\scanpage{27}
% L1, cancelled formula at left
\margin{$\deleted{M\xrightarrow{d^n}P^n_{A/\Lambda}(M)}$}
% L1–3
en correspondance biunivoque avec les éléments de \unsure{section|sections} d'ordre $n$ de $X\times_S Y$ \deleted{le long} sur $X\times_S X$ le long de la $=$ immersion \unsure{\ldots|{}}. \deleted{\unsure{$\Gamma_\varphi$}}

% L4–8
Soit $n'\geq n$, alors on a
\[
\Delta^{(\deleted{\unsure{n'|m}}n)}_{X/S}\longrightarrow\Delta^{(n')}_{X/S}
\]
et les $\Delta^{(n)}_{X/S}$ forment un \uline{système inductif} de schémas sur $X$, dont la limite est un \uline{schéma formel} sur $X^{\unsure{\to|{}}}$, noté $\Delta^{(\infty)}_{X/S}$.

% L9–12
On a évidemment
\[
\pi^n(Y/X/S)=\pi^{n'}(Y/X/S)\times_{\Delta^{(n')}_{X/S}}\Delta^{(n)}_{X/S}
\]
d'où aussitôt un $X$-hom naturel
\[
P^{n'}(Y/X/S)\xrightarrow{\varphi_{nn'}}p^n(Y/X/S).
\]
% L13–16
C'est un système \uline{transitif} de $X$-hom. et on a \deleted{$\unsure{\varphi}$}
\[
\left\{\begin{aligned}
\underline{\varphi_{\unsure{0|\circ|o}n}=\varepsilon^n}&\quad\text{d'où}\\
\underline{\underline{\varepsilon^{(n)}}}(Y/X/S)\circ\varphi_{nn'}&=\varepsilon^{(n')}\\
\varphi_{nn'}d^{(n')}&=d^{(n)}.
\end{aligned}\right.
\]
% archive sheet number: 26

